\section*{Hoofdstuk \ref{ch:g11:plate-motions} --- Platen en hun bewegingen}

\begin{solution}{exo:g11:plate-motions:1}
Een \omterm{def:g11:plate-motions:rigid-plate}{starre plaat} is een stuk \omterm{def:g11:lithosphere-asthenosphere:lithosphere}{lithosfeer} dat als één blok beweegt: de
afstanden tussen zijn punten veranderen niet, en zijn vervorming is samengebracht
in smalle zones langs zijn grenzen. De aardbevingen laten dat zien: bijna allemaal
vinden ze plaats langs smalle gordels, de plaatgrenzen, en heel weinig binnen de
platen.
\end{solution}

\begin{solution}{exo:g11:plate-motions:2}
\qty{72}{mm} per jaar is \qty{72}{km} per miljoen jaar. In 30 jaar beweegt het
station $72 \times 30 = \qty{2160}{mm}$, ongeveer \qty{2.2}{m}; in een miljoen jaar
\qty{72}{km}.
\end{solution}

\begin{solution}{exo:g11:plate-motions:3}
$v_E = 30/10 = \qty{3}{cm}$ per jaar en $v_N = 40/10 = \qty{4}{cm}$ per jaar. De
snelheid is $\sqrt{3^2 + 4^2} = \qty{5}{cm}$ per jaar. De hoek $\beta$ met het
noorden voldoet aan $\tan\beta = 3/4$, dus $\beta \approx \ang{37}$: het station
beweegt naar het noordoosten, \ang{37} ten oosten van het noorden.
\end{solution}

\begin{solution}{exo:g11:plate-motions:4}
Een \omterm{def:g11:plate-motions:hotspot-track}{hotspotspoor} is de rij vulkanen en onderzeese bergen die een plaat van een
hotspot wegvoert terwijl ze eroverheen beweegt, steeds ouder naarmate ze verder
van de actieve vulkaan liggen. De jongste vulkaan, de Kilauea, staat aan het
zuidoostelijke einde van de keten, op het eiland Hawaii; de ouderdommen nemen toe
naar het west-noordwesten. De plaat voert haar vulkanen van de hotspot weg in de
richting van haar beweging: de Pacifische plaat beweegt naar het
west-noordwesten.
\end{solution}

\begin{solution}{exo:g11:plate-motions:5}
$519/5.1 \approx \qty{102}{km}$ per miljoen jaar, dat wil zeggen \qty{102}{mm} per
jaar, ongeveer \qty{10}{cm} per jaar.
\end{solution}

\begin{solution}{exo:g11:plate-motions:6}
$v_E = 63/31 \approx \qty{20.3}{mm}$ per jaar en $v_N = 48/31 \approx
\qty{15.5}{mm}$ per jaar. De snelheid is $\sqrt{20.3^2 + 15.5^2} \approx
\qty{25.5}{mm}$ per jaar; $\tan\alpha = 15.5/20.3 \approx 0.76$ geeft $\alpha
\approx \ang{37}$: het station beweegt \ang{37} ten noorden van het oosten, naar het
noordoosten. De Mauna Kea beweegt $72/25.5 \approx 2.8$ keer sneller.
\end{solution}

\begin{solution}{exo:g11:plate-motions:7}
(1) In millimeter per jaar,
\[
  \vec v_{\mathrm{NA/EU}} = (-10.5 - 7.7,\ 20.6 - 16.3) = (-18.2, +4.3).
\] De grootte is $\sqrt{18.2^2 + 4.3^2} \approx
\qty{18.7}{mm}$ per jaar, en $\tan\alpha = 4.3/18.2 \approx 0.24$ geeft
\ang{13} ten noorden van het westen. (2) De twee platen bewegen met ongeveer
\qty{1.9}{cm} per jaar uit elkaar over de rug: dat is de \omterm{def:g11:drift-to-spreading:spreading-rate}{volledige
spreidingssnelheid} van de rug in IJsland. (3) Ze komt overeen met de \qty{2}{cm}
per jaar die voor de Noord-Atlantische Oceaan wordt genoemd.
\end{solution}

\begin{solution}{exo:g11:plate-motions:8}
(1) $v = \pi \times 0.651 \times 6371/180 \approx \qty{72.4}{km}$ per miljoen
jaar, ongeveer \qty{7.2}{cm} per jaar. (2) $72.4 \times \sin\ang{30} \approx
\qty{36.2}{km}$ per miljoen jaar, ongeveer \qty{3.6}{cm} per jaar. (3) Bij de
rotatiepool, als die op de plaat ligt. De snelheid verandert van punt tot
punt, van nul bij de pool tot haar grootste waarde op \ang{90} ervan: een
hele plaat heeft geen enkelvoudige snelheid, alleen een snelheid in elk punt (of
een gemiddelde over haar oppervlak).
\end{solution}

\begin{solution}{exo:g11:plate-motions:9}
(1) In millimeter per jaar,
\[
  \vec v_{\mathrm{IN/EU}} = (40.6 - 24.0,\ 38.6 + 3.3) = (+16.6, +41.9).
\] De grootte is $\sqrt{16.6^2 + 41.9^2} \approx
\qty{45}{mm}$ per jaar; de hoek $\beta$ met het noorden voldoet aan $\tan\beta =
16.6/41.9 \approx 0.40$, dus $\beta \approx \ang{22}$: India beweegt naar
Eurazië toe, \ang{22} ten oosten van het noorden. (2) \qty{45}{km} per miljoen
jaar, dus \qty{450}{km} in 10 miljoen jaar.
\end{solution}

\begin{solution}{exo:g11:plate-motions:10}
$(2432 - 1058)/(27.7 - 10.3) = 1374/17.4 \approx \qty{79}{km}$ per miljoen
jaar, ongeveer \qty{7.9}{cm} per jaar, iets minder dan \qty{9.1}{cm} per jaar.
De ouderdommen zijn elk onzeker met ongeveer een half miljoen jaar; het
ouderdomsverschil, $17.4 \pm 1.0$ miljoen jaar, geeft een snelheid tussen ongeveer 75 en
\qty{84}{km} per miljoen jaar. Tussen deze twee vulkanen lijkt de plaat iets
langzamer te zijn bewogen dan gemiddeld.
\end{solution}

\begin{solution}{exo:g11:plate-motions:11}
(1) $\vec v_{A/C} = \vec v_{A/B} + \vec v_{B/C}$: \qty{3}{cm} per jaar naar het oosten plus
\qty{4}{cm} per jaar naar het noorden. De twee pijlen staan loodrecht op elkaar, en
de derde zijde van de driehoek heeft de lengte $\sqrt{3^2 + 4^2} = \qty{5}{cm}$ per jaar,
in de richting \ang{37} ten oosten van het noorden ($\tan\beta = 3/4$). (2) $\vec v_{C/A}
= -\vec v_{A/C}$: \qty{5}{cm} per jaar in de tegenovergestelde richting, \ang{37} ten
westen van het zuiden.
\end{solution}

\begin{solution}{exo:g11:plate-motions:12}
(1) $54/0.375 = \qty{144}{km}$ per miljoen jaar, \qty{14.4}{cm} per jaar. (2)
Van $29/0.375 \approx 77$ tot $79/0.375 \approx \qty{211}{km}$ per miljoen
jaar: tussen ongeveer 8 en \qty{21}{cm} per jaar. (3) Dicht bij de hotspot is de
afstand niet groter dan de afmeting van een vulkaan of van het gebied waarover
de hotspot vulkanen voedt, en de ouderdom niet groter dan de tijd die een vulkaan
nodig heeft om te groeien: beide zijn relatief gezien heel onzeker. Ver van de
hotspot is een fout van \qty{25}{km} op \qty{2400}{km}, of van een half miljoen jaar
op 28 miljoen, maar een of twee procent.
\end{solution}

\begin{solution}{exo:g11:plate-motions:13}
(1) Tien keer de spreiding is \qty{50}{mm}: $50/72 \approx 0.7$ jaar, ongeveer
acht maanden. (2) $50/5 = 10$ jaar. (3) Over miljoenen jaren bewegen de platen
tientallen tot honderden kilometers, afstanden die de magnetische strepen
vastleggen en die op zee makkelijk te meten waren; plaatsbepaling tot op een paar
millimeter, die een beweging in maanden zichtbaar maakt, kwam veel later.
\end{solution}

\begin{solution}{exo:g11:plate-motions:14}
(1) In millimeter per jaar,
\[
  \vec v_{\mathrm{SA/AF}} = (-4.9 - 27.1,\ 11.0 - 17.9) = (-32.0, -6.9).
\] De grootte is $\sqrt{32.0^2 + 6.9^2} \approx
\qty{32.7}{mm}$ per jaar, en $\tan\alpha = 6.9/32.0 \approx 0.22$ geeft
\ang{12} ten zuiden van het westen. (2) De \omterm{def:g11:drift-to-spreading:spreading-rate}{volledige spreidingssnelheid} is ongeveer
\qty{3.3}{cm} per jaar, de halve ongeveer \qty{1.6}{cm} per jaar. (3) De halve
snelheid van de laatste paar miljoen jaar, \qty{1.6}{cm} per jaar, is iets kleiner
dan de \qty{1.9}{cm} per jaar gemiddeld over 66 miljoen jaar: de Zuid-Atlantische
Oceaan is recent misschien iets langzamer opengegaan dan gemiddeld, of het verschil
kan komen van de onzekerheden in de ouderdom van de sedimenten. Beide metingen zijn
het erover eens dat de rug langzaam is.
\end{solution}

\begin{solution}{exo:g11:plate-motions:15}
(1) Een rechte lijn die met \qty{20}{mm} per jaar stijgt, een plotselinge daling van
\qty{50}{mm} na vijf jaar, en dan een lijn die weer met dezelfde helling stijgt.
(2) Het station eindigt $10 \times 20 - 50 = \qty{150}{mm}$ ten oosten van zijn begin:
$150/10 = \qty{15}{mm}$ per jaar, een kwart te weinig. (3) Elke helft geeft
\qty{20}{mm} per jaar. (4) De sprong is een afzonderlijke, plotselinge beweging die de
aardbeving veroorzaakt: de geodeet meet de grootte ervan uit de posities vlak ervoor
en vlak erna, haalt haar uit de reeks (of past de twee delen aan met dezelfde
helling en een stap ertussen), en houdt \qty{20}{mm} per jaar aan als de
snelheid van het station.
\end{solution}

\begin{solution}{pb:g11:plate-motions:1}
\textbf{1.} De Mauna Kea (\num{0.375} miljoen jaar) is de jongste, Suiko
(midden) (\num{64.7} miljoen jaar) de oudste. De ouderdommen nemen toe met de
afstand tot de Kilauea.

\textbf{2.} $\qty{1}{km} = 10^6~\unit{mm}$ en een miljoen jaar is $10^6$
jaar: \qty{1}{km} per miljoen jaar is $10^6~\unit{mm}/10^6$ jaar,
\qty{1}{mm} per jaar.

\textbf{3.} $519/5.1 \approx \qty{102}{mm}$ per jaar.

\textbf{4.} Necker $1058/10.3 \approx 102.7$; Laysan $1818/19.9 \approx
91.4$; Midway $2432/27.7 \approx 87.8$; Colahan $3128/38.6 \approx 81.0$;
Abbott $3280/38.7 \approx 84.8$ (millimeter per jaar).

\textbf{5.} De punten (54; 0.375), (519; 5.1), (780; 7.2), (1058; 10.3),
(1818; 19.9), (2432; 27.7), (3128; 38.6), (3280; 38.7) liggen dicht bij een rechte
lijn door de oorsprong, met enige spreiding (tegenover een lijn van \qty{91}{km} per
miljoen jaar liggen Kauai, Nihoa en Necker er iets onder, Laysan erop,
Colahan en Abbott er iets boven).

\textbf{6.} Elke vulkaan is boven de hotspot gebouwd en is sindsdien met de
plaat meebewogen. In de tijd $t$ is de plaat $v\,t$ in een constante
richting bewogen, dus de vulkaan ligt op $d = v\,t$ van het punt dat nu boven de
hotspot ligt: $t = d/v$.

\textbf{7.} $(102.7 + 91.4 + 87.8 + 81.0 + 84.8)/5 \approx \qty{89.5}{mm}$ per
jaar.

\textbf{8.} 89.5 tegenover \qty{91}{km} per miljoen jaar: ze verschillen
minder dan \qty{2}{\percent}.

\textbf{9.} $519/91 \approx 5.7$ miljoen jaar, tegenover een ouderdom van \num{5.1}
miljoen jaar: dichtbij, waarbij Kauai iets jonger is dan de gemiddelde lijn
voorspelt.

\textbf{10.} De vulkanen zijn steeds ouder naar het west-noordwesten:
de plaat voert ze in die richting mee. De Pacifische plaat beweegt naar het
west-noordwesten over de hotspot.

\textbf{11.} $(4860 - 3493)/(64.7 - 42.4) = 1367/22.3 \approx \qty{61}{km}$
per miljoen jaar, ongeveer \qty{6.1}{cm} per jaar.

\textbf{12.} $65/91 \approx 0.71$: langs de Keizerketen bewoog de plaat
ongeveer \qty{30}{\percent} langzamer dan langs de Hawaiiketen.

\textbf{13.} $3668/50 \approx \qty{73}{km}$ per miljoen jaar, ongeveer
\qty{7.3}{cm} per jaar, minder dan \qty{91}{km} per miljoen jaar. De lijn door de oorsprong
gebruikt de ouderdommen van 1987, die het einde van de Hawaiiketen op 39 tot 43
miljoen jaar leggen; de nieuwere ouderdom van ongeveer 50 miljoen jaar voor het begin
van de knik spreidt dezelfde afstand over een langere tijd. Beide kloppen als de plaat
in het vroege deel van de Hawaiiketen langzamer bewoog en later sneller:
de snelheid was niet constant.

\textbf{14.} Een verandering in de bewegingsrichting van de Pacifische plaat, ongeveer
50 miljoen jaar geleden.

\textbf{15.} De richting van de keten is de richting waarin de plaat
bewoog; de spreiding van de ouderdommen erlangs is haar snelheid. Een plaat kan
versnellen of vertragen zonder te draaien, of draaien zonder van snelheid te
veranderen, dus de twee zijn onafhankelijk. Snelheidsveranderingen zonder
richtingsverandering langs de Keizerketen laten zien dat haar snelheidsregistratie
en de knik afzonderlijke gebeurtenissen zijn: de knik is niet ontstaan doordat de
hotspot vertraagde.

\textbf{16.} $\sqrt{63.2^2 + 34.5^2} \approx \qty{72.0}{mm}$ per jaar, en
$\tan\alpha = 34.5/63.2 \approx 0.55$ geeft \ang{29} ten noorden van het westen.

\textbf{17.} West-noordwest is \ang{22.5} ten noorden van het westen, noordwest
\ang{45}: de richting van het station, \ang{29} ten noorden van het westen, ligt
tussen de twee, dicht bij de richting van de keten. Ze is ermee in overeenstemming.

\textbf{18.} De snelheid van het station is gemeten ten opzichte van een stelsel
waarin de \omterm{def:g11:lithosphere-asthenosphere:lithosphere}{lithosfeer} als geheel niet draait, terwijl de keten de beweging ten
opzichte van de hotspot geeft: als de twee stelsels ten opzichte van elkaar
bewegen, verschillen de snelheden. En het station meet de laatste 30 jaar,
de keten de laatste 40 miljoen jaar, waarin de snelheid veranderd kan zijn.
De spreiding van de ouderdommen laat ook een onzekerheid van een paar millimeter per
jaar op de snelheid van de keten.

\textbf{19.} $100/91 \approx 1.1$ miljoen jaar.

\textbf{20.} \textbf{De Pacifische plaat is de laatste 40 miljoen jaar met
ongeveer \qty{9}{cm} per jaar over de hotspot van Hawaii bewogen}, naar het
west-noordwesten; vóór de knik bewoog ze langzamer, en in een andere
richting.
\end{solution}
